\documentclass[a4paper,11pt]{article}
\usepackage[utf8]{inputenc}
\usepackage[T1]{fontenc}
\usepackage{lmodern}
\usepackage[french]{babel}
\usepackage{amsmath,amssymb}
\usepackage{booktabs,array}
\usepackage{siunitx}
\sisetup{locale=FR}
\usepackage[a4paper,margin=2.5cm]{geometry}
\usepackage{xcolor}
\usepackage[colorlinks=true,linkcolor=blue!50!black,urlcolor=blue!50!black]{hyperref}

\title{\textbf{Optimisation de l'amortissement des articulations}\\[2mm]
\large Démarche, critère et résultats}
\author{Hugo TESTANIERE}
\date{}

\begin{document}
\maketitle
\tableofcontents

\section{Objectif}
Les articulations du convoi reçoivent un couple amortisseur $M_k=c_k\,(\dot\theta_{k+1}-\dot\theta_k)$ (partie commande du rapport). Sans amortissement, le convoi présente une divergence de la tête et un serpentement instables dès quelques km/h. Avec un amortissement trop faible il reste instable ; trop fort, il devient inefficace (couples saturés) et réduit la manœuvrabilité.

L'optimisation cherche, pour chaque nombre de portions $N=3,\dots,10$ et chaque vitesse $v\in\{10,20,30,40\}$\,km/h, le \emph{coefficient d'amortissement moyen} $c$ et la \emph{répartition} de cet amortissement le long du convoi qui donnent le meilleur suivi de trajectoire. Les résultats alimentent la page de simulation : changer $N$ ou la vitesse règle automatiquement $c$ et la répartition.

\section{Variables optimisées}
\begin{itemize}
  \item $c$ : amortissement \emph{moyen} des $n=N-1$ articulations (\si{N.m.s/rad}).
  \item $s\in[-0{,}9;\,0{,}9]$ : \emph{pente} de la répartition linéaire. Pour l'articulation $k$ ($k=0,\dots,n-1$ de l'avant vers l'arrière), avec $x_k=k/(n-1)$,
  \begin{equation}
    g_k=\max\big(0,\;1+s\,(1-2x_k)\big),\qquad
    c_k=c\;\frac{g_k}{\frac1n\sum_j g_j}.
  \end{equation}
  La moyenne des $c_k$ vaut toujours $c$ : on compare des répartitions à amortissement total constant. $s>0$ concentre l'amortissement vers l'avant, $s<0$ vers l'arrière, $s=0$ est la répartition uniforme.
\end{itemize}

\section{Ce qui est fixé}
\begin{center}\small
\begin{tabular}{@{}ll@{}}\toprule
Paramètre & Valeur\\\midrule
Modèle & Lagrange non linéaire, intégration Runge--Kutta d'ordre 4\\
Pneus & rigidité $K_t=\SI{13600}{N/rad}$, saturation en $\tanh$, $\mu=0{,}8$\\
Articulation 1 & bloquée (réglage par défaut de la page)\\
Couples & amortisseur idéal, saturés à $\pm\SI{2500}{N.m}$ (sans dynamique des vérins)\\
Pilote & loi de Stanley sur l'essieu de tête, gain $1$, braquage limité à $\pm35^\circ$\\
Masses, longueurs & valeurs par défaut (voir page de simulation)\\\bottomrule
\end{tabular}
\end{center}

\section{Manœuvres de test}
Quatre manœuvres, chacune simulée à la vitesse étudiée :
\begin{enumerate}
  \item \textbf{Test de l'élan} : couloirs de plots de l'ISO 3888-2 pour un véhicule de largeur $b=\SI{1,5}{m}$ (entrée \SI{12}{m} de large $1{,}1\,b+0{,}25$, transition \SI{13,5}{m}, évitement \SI{11}{m} de large $b+1$ décalé de \SI{1}{m} par rapport au bord du premier couloir, retour \SI{12,5}{m}, sortie \SI{12}{m}), 18 plots, trajectoire de référence suivant l'axe des couloirs (polynôme de degré 5 dans les transitions).
  \item \textbf{Slalom} : amplitude \SI{1,5}{m}, plots espacés de \SI{12}{m}, $5$ plots.
  \item \textbf{Virage} de $90^\circ$ à gauche, rayon \SI{12}{m}.
  \item \textbf{Freinage en virage} : virage de $90^\circ$, rayon \SI{20}{m}, décélération \SI{3}{m/s^2} à partir de $30\,\%$ de l'arc, jusqu'à \SI{5}{km/h}.
\end{enumerate}

\section{Critère}
Pour une manœuvre $m$ on mesure sur la simulation :
\begin{itemize}
  \item $e_m$ : erreur maximale du centre d'essieu de la dernière portion par rapport à la trajectoire (m) ;
  \item $d_m$ : \emph{déport} maximal, écart latéral entre la queue et la trajectoire réellement suivie par la tête (m) ;
  \item $p_m$ : nombre de plots touchés (test de l'élan seulement) ;
  \item $\mathrm{RWA}_m$ : \emph{rearward amplification}, rapport de l'accélération latérale maximale de la dernière portion à celle de la première (accélérations au centre de gravité, norme ISO~14791), borné à $5$.
\end{itemize}
Le coût d'une manœuvre est
\begin{equation}
  J_m=\begin{cases}
    e_m+0{,}5\,d_m+0{,}15\,\mathrm{RWA}_m+0{,}1\,p_m & \text{si le convoi suit la trajectoire,}\\
    10 & \text{si le convoi diverge (perte de contrôle).}
  \end{cases}
\end{equation}
Le coût d'une configuration $(c,s)$ combine le pire cas et la moyenne des manœuvres, avec une très légère pénalité sur $c$ pour départager les cas équivalents :
\begin{equation}
  J(c,s)=\tfrac12\max_m J_m+\tfrac12\,\overline{J_m}+0{,}01\,\ln c .
\end{equation}
Les poids ($0{,}5$ pour le déport, $0{,}15$ pour le RWA) sont des choix de l'auteur : un point de RWA en plus vaut environ \SI{15}{cm} d'erreur de suivi.

\section{Algorithme}
Pour chaque couple $(N,v)$ :
\begin{enumerate}
  \item \textbf{Grille grossière} : $c\in\{1000;\,2500;\,6000;\,15000;\,40000\}$ et $s\in\{-0{,}6;\,-0{,}3;\,0;\,0{,}3;\,0{,}6\}$, soit $25$ configurations, chacune simulée sur les $4$ manœuvres.
  \item \textbf{Raffinement local} (deux itérations) autour du meilleur point : on teste $c\times\{1/f,1,f\}$ et $s+\{-\delta,0,\delta\}$ avec $f=1{,}6$ puis $\sqrt{1{,}6}$, et $\delta=0{,}15$ puis $0{,}075$.
\end{enumerate}
C'est une recherche locale, pas un optimiseur global : l'optimum trouvé n'est pas garanti global. Environ $35$ configurations $\times\,4$ manœuvres sont simulées par couple $(N,v)$.

\section{Résultats}
Les tableaux donnent le réglage retenu, le coût $J$, le RWA moyen sur les manœuvres qui ne divergent pas, et les manœuvres où le convoi diverge malgré le meilleur réglage trouvé.
\begin{table}[ht]\centering\small
\begin{tabular}{@{}r r r r r l@{}}\toprule
$N$ & $c$ (N$\cdot$m$\cdot$s/rad) & pente $s$ & $J$ & RWA moyen & manoeuvres divergentes\\\midrule
3 & 494 & $-0{,}82$ & 0{,}39 & 1{,}03 & --\\
4 & 1563 & $-0{,}82$ & 0{,}38 & 1{,}01 & --\\
5 & 4743 & $+0{,}37$ & 0{,}38 & 0{,}96 & --\\
6 & 7589 & $+0{,}45$ & 0{,}37 & 0{,}88 & --\\
7 & 4743 & $+0{,}82$ & 0{,}36 & 0{,}81 & --\\
8 & 4743 & $+0{,}60$ & 0{,}36 & 0{,}81 & --\\
9 & 6000 & $+0{,}67$ & 0{,}37 & 0{,}81 & --\\
10 & 7589 & $+0{,}52$ & 0{,}39 & 0{,}82 & --\\
\bottomrule\end{tabular}
\caption{Réglage optimal à \SI{10}{km/h}.}\label{tab:opt10}\end{table}

\begin{table}[ht]\centering\small
\begin{tabular}{@{}r r r r r l@{}}\toprule
$N$ & $c$ (N$\cdot$m$\cdot$s/rad) & pente $s$ & $J$ & RWA moyen & manoeuvres divergentes\\\midrule
3 & 494 & $+0{,}82$ & 0{,}33 & 1{,}11 & --\\
4 & 2500 & $+0{,}75$ & 0{,}41 & 1{,}08 & --\\
5 & 1976 & $+0{,}82$ & 0{,}50 & 1{,}11 & --\\
6 & 1976 & $+0{,}82$ & 0{,}63 & 1{,}11 & --\\
7 & 1563 & $+0{,}75$ & 0{,}74 & 1{,}20 & --\\
8 & 2965 & $+0{,}82$ & 0{,}96 & 1{,}04 & --\\
9 & 2965 & $+0{,}82$ & 1{,}14 & 1{,}07 & --\\
10 & 2965 & $+0{,}82$ & 1{,}41 & 1{,}15 & --\\
\bottomrule\end{tabular}
\caption{Réglage optimal à \SI{20}{km/h}.}\label{tab:opt20}\end{table}

\begin{table}[ht]\centering\small
\begin{tabular}{@{}r r r r r l@{}}\toprule
$N$ & $c$ (N$\cdot$m$\cdot$s/rad) & pente $s$ & $J$ & RWA moyen & manoeuvres divergentes\\\midrule
3 & 15000 & $-0{,}60$ & 1{,}56 & 1{,}10 & --\\
4 & 15000 & $+0{,}60$ & 2{,}78 & 0{,}97 & --\\
5 & 6000 & $+0{,}45$ & 7{,}10 & 0{,}97 & slalom\\
6 & 15000 & $+0{,}45$ & 7{,}50 & 1{,}02 & slalom\\
7 & 3750 & $+0{,}22$ & 7{,}74 & 1{,}26 & slalom\\
8 & 6000 & $+0{,}60$ & 7{,}84 & 1{,}20 & slalom\\
9 & 7589 & $+0{,}22$ & 8{,}36 & 0{,}86 & test de l'élan, slalom\\
10 & 3162 & $+0{,}82$ & 8{,}94 & 0{,}22 & test de l'élan, slalom, virage 90°\\
\bottomrule\end{tabular}
\caption{Réglage optimal à \SI{30}{km/h}.}\label{tab:opt30}\end{table}

\begin{table}[ht]\centering\small
\begin{tabular}{@{}r r r r r l@{}}\toprule
$N$ & $c$ (N$\cdot$m$\cdot$s/rad) & pente $s$ & $J$ & RWA moyen & manoeuvres divergentes\\\midrule
3 & 24000 & $-0{,}37$ & 5{,}58 & 1{,}22 & --\\
4 & 24000 & $-0{,}07$ & 7{,}46 & 1{,}12 & virage 90°\\
5 & 15000 & $-0{,}82$ & 8{,}32 & 0{,}89 & test de l'élan, virage 90°\\
6 & 7589 & $+0{,}53$ & 9{,}13 & 0{,}77 & test de l'élan, slalom, virage 90°\\
7 & 7589 & $+0{,}67$ & 9{,}15 & 0{,}69 & test de l'élan, slalom, virage 90°\\
8 & 7589 & $+0{,}82$ & 9{,}16 & 0{,}59 & test de l'élan, slalom, virage 90°\\
9 & 7589 & $+0{,}82$ & 9{,}16 & 0{,}51 & test de l'élan, slalom, virage 90°\\
10 & 7589 & $+0{,}82$ & 9{,}17 & 0{,}42 & test de l'élan, slalom, virage 90°\\
\bottomrule\end{tabular}
\caption{Réglage optimal à \SI{40}{km/h}.}\label{tab:opt40}\end{table}


\paragraph{Lecture.}
\begin{itemize}
  \item À faible vitesse le minimum est plat : plusieurs couples $(c,s)$ donnent des résultats voisins, le réglage retenu est peu discriminant.
  \item À mesure que la vitesse et le nombre de portions augmentent, le coût croît et des divergences apparaissent : au-delà de \SI{30}{km/h}, aucun réglage testé ne stabilise les convois longs dans ces manœuvres.
  \item Test de l'élan (couloirs ISO~3888-2, $18$ plots) avec le réglage optimal : aucun plot touché à \SI{10}{km/h}, de $0$ à $2$ plots à \SI{20}{km/h}, de $1$ à $10$ à \SI{30}{km/h} et de $2$ à $7$ à \SI{40}{km/h} selon $N$ : le convoi suit la trajectoire mais ne tient pas dans les couloirs aux vitesses élevées.
  \item Lorsque des manœuvres divergent, les valeurs affichées sont le « moins mauvais » réglage testé : elles ne stabilisent pas le convoi.
\end{itemize}

\section{Utilisation dans la page de simulation}
La table est intégrée à la page (\texttt{doc/animation/simulation\_convoi.html}). La case « c et répartition optimisés selon N et la vitesse » règle $c$ et la pente quand on modifie $N$ ou la vitesse. Entre deux vitesses de la table, $c$ est interpolé en logarithme et la pente linéairement. Un avertissement s'affiche lorsque le réglage correspond à une configuration qui diverge. Modifier $c$ ou la répartition à la main désactive la case.

\section{Limites}
\begin{itemize}
  \item Résultats valables pour les conditions et manœuvres du tableau ci-dessus ; non garantis avec une autre adhérence, d'autres masses, d'autres amplitudes de manœuvre, sans saturation de pneus ou avec la dynamique des vérins.
  \item Famille de répartitions restreinte à la loi linéaire ; le gain du pilote, les gains du vérin et la saturation des couples ne sont pas optimisés.
  \item Le modèle de pneu ne prend pas en compte l'adhérence combinée ni le transfert de charge.
  \item Le modèle du convoi n'a pas été comparé numériquement au code MATLAB d'origine.
\end{itemize}

\section{Relancer le calcul}
Les fichiers sont dans \texttt{doc/optimisation/} :
\begin{itemize}
  \item \texttt{opt\_rwa.html} : programme d'optimisation (à ouvrir dans Chrome avec \texttt{\#V=20\&N=3-10}) ; la progression sort dans la console du navigateur ;
  \item \texttt{logs/} : journaux de calcul ;
  \item \texttt{generer.py} : fusionne les journaux, écrit \texttt{resultats.json} (lu par la page) et \texttt{tableaux.tex} (ce document) ;
  \item \texttt{v1\_sans\_RWA/} : première version de l'optimisation, sans le terme de RWA dans le critère.
\end{itemize}
Après un nouveau calcul : \texttt{python3 generer.py}, puis \texttt{python3 construire.py} dans \texttt{doc/animation/}.

\end{document}
